% !TeX program = pdflatex
% Based on ../poster/main.tex: A1 portrait, tikzposter, two columns, blue on white.
% Scientific source: ../../artifacts_article/icomp_artifacts.pdf (19 Sep 2026).
\documentclass[17pt,a1paper,portrait,margin=15mm,colspace=10mm,blockverticalspace=5mm]{tikzposter}
\usepackage[T1]{fontenc}
\usepackage[utf8]{inputenc}
\usepackage[english]{babel}
\usepackage{lmodern}
\usepackage{amsmath,amssymb}
\usepackage{graphicx,booktabs,array,tabularx,ragged2e}
\usepackage[hidelinks]{hyperref}
\usepackage{qrcode}
\tikzposterlatexaffectionproofoff
\usetheme{Default}
\renewcommand{\familydefault}{\sfdefault}
\definecolor{Blue}{HTML}{004488}
\definecolor{Ink}{HTML}{172B40}
\definecolor{Rose}{HTML}{BB5566}
\definecolor{Gold}{HTML}{997700}
\definecolor{Muted}{HTML}{52606D}
\definecolorstyle{AcademicStyle}{}{
 \colorlet{backgroundcolor}{white}\colorlet{framecolor}{white}
 \colorlet{titlefgcolor}{Blue}\colorlet{titlebgcolor}{white}
 \colorlet{blocktitlebgcolor}{Blue}\colorlet{blocktitlefgcolor}{Blue}
 \colorlet{blockbodybgcolor}{white}\colorlet{blockbodyfgcolor}{Ink}
}
\usecolorstyle{AcademicStyle}
\defineblockstyle{PosterMinimal}{
 titlewidthscale=1,bodywidthscale=1,titleleft,
 titleoffsetx=0pt,titleoffsety=0pt,bodyoffsetx=0pt,bodyoffsety=0pt,
 bodyverticalshift=0pt,roundedcorners=0,linewidth=1.2pt,
 titleinnersep=3mm,bodyinnersep=3mm
}{
 \ifBlockHasTitle
 \draw[color=Blue,line width=1.2pt] (blocktitle.south west)--(blocktitle.south east);
 \fi
}
\useblockstyle{PosterMinimal}
\newcommand{\pblock}[2]{\block{\fontsize{26}{30}\selectfont #1}{\fontsize{20}{25}\selectfont\RaggedRight\setlength{\parskip}{7pt}#2}}
\newcommand{\pcaption}[1]{{\fontsize{16}{20}\selectfont\color{Muted}#1\par}}
\newcommand{\plot}[1]{\par\vspace{2mm}\centerline{\includegraphics[width=\linewidth]{figures/#1.pdf}}\vspace{1mm}}
\newcommand{\emphline}[1]{{\color{Blue}\bfseries #1}}
\newcommand{\dact}{d_{\mathrm{act}}}
\newcommand{\dmg}{\hat d_{\mathrm{MG}}}
\DeclareMathOperator{\PR}{PR}
\newenvironment{posteritems}{\begin{itemize}\setlength{\itemsep}{3pt}\setlength{\parsep}{0pt}\setlength{\topsep}{3pt}}{\end{itemize}}
\title{Estimating Active Dimension of Training Dynamics}
\author{Nikolay Karlov\textsuperscript{1}, Alexey Kravatskiy\textsuperscript{2}}
\institute{\textsuperscript{1}Moscow Institute of Physics and Technology \quad \textsuperscript{2}AIRI}
\settitle{
 {\fontsize{19}{23}\selectfont\bfseries\color{Blue} ASCOMP\enspace /\enspace ICOMP 2026\hfill MIPT\enspace /\enspace AIRI\par}
 \vspace{6mm}
 {\fontsize{48}{55}\selectfont\bfseries\color{Blue}Estimating Active Dimension\\of Training Dynamics from One Scalar Log\par}
 \vspace{3mm}
 {\fontsize{27}{33}\selectfont\color{Ink}with an Application to Grokking\par}
 \vspace{5mm}
 {\fontsize{21}{27}\selectfont\color{Ink}\textbf{Nikolay Karlov\textsuperscript{1}\quad Alexey Kravatskiy\textsuperscript{2}}\par}
 \vspace{2mm}
 {\fontsize{19}{23}\selectfont\color{Ink}\textsuperscript{1}Moscow Institute of Physics and Technology\qquad\textsuperscript{2}AIRI\par}
}
\begin{document}
\maketitle[width=564mm,innersep=3mm,titletotopverticalspace=0mm,titletoblockverticalspace=5mm]
\block[bodyinnersep=4mm]{}{
 \colorbox{Blue!7}{\parbox{\dimexpr\linewidth-2\fboxsep\relax}{\fontsize{23}{29}\selectfont\color{Blue}\textbf{One log can reveal low-dimensional recurrent dynamics.}\par
 \fontsize{21}{27}\selectfont\color{Ink}On the grokking runs studied here, it gives a transition signal; stored trajectories reveal a reversible change in covariance effective rank.}}
}
\begin{columns}
\column{0.5}
\pblock{01\quad Dimensions and covariance effective rank}{
Let $\theta=\theta_0+V^{\top}c$, where $V\in\mathbb R^{k\times P}$, $VV^{\top}=I_k$ and $c\in\mathbb R^k$. Let $f(c)$ be the outputs on a fixed probe set.

Let $\mu_R$ be the occupation measure induced by an ergodic invariant measure of regime $R$, assumed exactly dimensional. Write $C_{\mathcal W}=\operatorname{Cov}_{t\in\mathcal W}(c_t)$ for a trajectory window $\mathcal W$.

{\fontsize{18}{23}\selectfont\renewcommand{\arraystretch}{1.25}
\begin{tabularx}{\linewidth}{@{}>{\bfseries\RaggedRight\arraybackslash}p{0.31\linewidth}X@{}}
\toprule
Quantity & \textbf{Definition}\\
\midrule
Available dimension & $d_{\mathrm{avail}}=\dim\operatorname{im}(V^{\top})=k$.\\
Functional dimension & $d_{\mathrm{func}}(c)=\operatorname{rank}\!\left(\dfrac{\partial f(c)}{\partial c}\right)$.\\
Active dimension & $\dact(R)=\displaystyle\lim_{\varepsilon\downarrow0}\dfrac{\log\mu_R(B(c,\varepsilon))}{\log\varepsilon}$, $\mu_R$-a.e.\\
Covariance effective rank & $\PR^{\mathrm{pos}}=\dfrac{(\operatorname{tr}C_{\mathcal W})^2}{\operatorname{tr}(C_{\mathcal W}^2)}$, \quad $C_{\mathcal W}\ne0$.\\
\bottomrule
\end{tabularx}}
\pcaption{$B(c,\varepsilon)$ is a Euclidean ball. For $\PR^{\mathrm{det}}$, replace $c_t$ by its residual about the window's least-squares line. Finite samples resolve finite-scale exponents, not the limit defining $\dact$.}

}
\pblock{02\quad Delay embedding and dimension estimation}{
\plot{fig_method}
\pcaption{One scalar record, its delay reconstruction, and a neighbourhood. The loop is visited repeatedly.}
\[
 y_t=(x_t,x_{t-\tau},\ldots,x_{t-(E-1)\tau})\in\mathbb R^E,
\]
\[
 \dmg^{-1}=\frac{1}{N(m-1)}\sum_{i=1}^{N}\sum_{j=1}^{m-1}
 \log\frac{r_m^{(i)}}{r_j^{(i)}}.
\]
Pool the nearest-neighbour likelihood over $N$ delay vectors, each with $m$ neighbours, before inversion. $r_j^{(i)}$ is the $j$th neighbour distance; the \textbf{Theiler exclusion} requires time separation greater than $W_T$.

\emphline{A dimensional interpretation needs} a valid delay embedding, one regime per window, enough returns after exclusion, and a stable scaling range.
}
\pblock{03\quad Recovery accuracy across dynamical regimes}{
A linear head on a frozen image network is forced by $r$ independent quasiperiodic phases. This supplies known active dimension. Settings are frozen before testing held-out ranks and seeds.
\plot{fig_regimes}
\pcaption{Left: recovery by regime; markers are medians over seeds and observers, bands are interquartile ranges. Right: stability under doubled embedding, at $r=2$ (open) and $r=6$ (filled).}
\begin{posteritems}
\item \textcolor{Blue}{\textbf{Recurrent:}} about one-component error up to roughly eight components at the main configuration.
\item \textcolor{Rose}{\textbf{Stochastic:}} noise shifts the value and makes it depend on embedding settings.
\item \textcolor{Gold}{\textbf{Transient:}} temporal exclusion can produce a large spurious value.
\end{posteritems}
\pcaption{Only one learned-system benchmark passes every prescribed check. All validation systems are externally forced.}
}
\column{0.5}
\pblock{04\quad Scalar-log diagnostics in grokking}{
\textbf{Grokking} is delayed generalisation after memorisation. We examine a mini-batch transformer and an unregularised full-batch perceptron.
\plot{fig_map}
\pcaption{Each point is a log. The blue band records accurate controlled cases; it is not a validated decision boundary.}
The \textbf{identifiability ratio}
\[
 \rho_{\mathrm{ident}}=\frac{\dmg(2E_{\max})}{\dmg(E_{\max})}
\]
tests sensitivity to embedding dimension. Trend crossings measure non-monotonicity of the log.

\emphline{All evaluated grokking logs lie outside the reference region.} Their scalar levels are not interpreted as active dimension. These diagnostics test neither recurrence itself nor the scaling range.
}
\pblock{05\quad Covariance effective rank near generalisation}{
We record parameter and probe-logit trajectories through 1024-dimensional CountSketch projections. In each window, detrended covariance eigenvalues $\lambda_i$ give
\[
 \PR^{\mathrm{det}}=\frac{(\sum_i\lambda_i)^2}{\sum_i\lambda_i^2}.
\]
\plot{fig_dip}
\pcaption{Four regularised runs (rose): three modular-addition seeds and one $S_5$ task. Two runs without weight decay (gold) are aligned to paired transition times. Thick rose curves average four runs; shading is hash-family spread, not uncertainty across runs.}
\begin{posteritems}
\item Covariance effective rank falls near generalisation and later re-expands in parameter and function space.
\item A window matched to the transition gives a correlated drop in the parameter-norm log estimate.
\end{posteritems}
\emphline{Effective rank can fall while active dimension stays fixed.} The scalar drop describes a transition retrospectively; it does not establish early prediction.
}
\end{columns}
\pblock{06\quad Estimation limits and experimental confounders}{
\begin{minipage}[t]{0.485\linewidth}
\begin{posteritems}
\item Recovery depends on the lag, observer and window. The tested resolution ceiling is low.
\item The grokking comparison has four regularised runs against two without weight decay. Regularisation and outcome are confounded.
\end{posteritems}
\end{minipage}\hfill
\begin{minipage}[t]{0.485\linewidth}
\begin{posteritems}
\item One memorising run falls just as deeply elsewhere. Depth alone does not distinguish generalisation.
\item The scalar effect disappears after linear detrending; roughness changes more strongly. The trajectory result does not reproduce in the full-batch setting.
\end{posteritems}
\end{minipage}\par\vspace{2mm}
\emphline{Interpret the geometry only after checking the regime.} For these grokking runs, the evidence supports a reversible redistribution of trajectory variance.
}
\block[bodyinnersep=3mm]{}{
\color{Blue}\rule{\linewidth}{1pt}\par\vspace{2mm}
\begin{minipage}[c]{0.80\linewidth}
{\fontsize{16}{20}\selectfont\color{Ink}
\textbf{Method references:} Takens (1981); Sauer et al. (1991); Levina \& Bickel (2004); MacKay \& Ghahramani (2005).\par
\vspace{2mm}\textbf{Contact:} \href{mailto:karlov.na@phystech.edu}{karlov.na@phystech.edu}\quad\textbf{Paper \& code:} intsystems/2026-Project-202}
\end{minipage}\hfill
\begin{minipage}[c]{0.16\linewidth}\centering
\qrcode[height=22mm]{https://github.com/intsystems/2026-Project-202}\par
{\fontsize{14}{17}\selectfont\color{Muted}Paper \& code}
\end{minipage}
}
\end{document}

